
\documentclass[UTF8,11pt,a4paper]{ctexart}
\usepackage[a4paper,left=23mm,right=23mm,top=24mm,bottom=24mm]{geometry}
\usepackage{fontspec}
\setmainfont{DejaVu Sans}
\setCJKmainfont{Noto Serif CJK SC}
\usepackage{amsmath,amssymb}
\usepackage{xcolor}
\usepackage[colorlinks=true,linkcolor=blue!45!black,unicode=true]{hyperref}
\usepackage{enumitem}
\usepackage{needspace}
\usepackage{fancyhdr}
\setlength{\parindent}{2em}
\setlength{\parskip}{0.25em}
\setlist[enumerate]{leftmargin=2.6em,itemsep=0.4em,topsep=0.5em}
\pagestyle{fancy}
\fancyhf{}
\fancyhead[L]{\small 三千弱水 · 大学化学高阶综合题库}
\fancyhead[R]{\small 待专家审校稿}
\fancyfoot[C]{\thepage}
\setlength{\headheight}{15pt}
\sloppy

\begin{document}
\begin{center}
{\LARGE\bfseries 大学化学高阶综合题 · 第8辑 (5题)}\\[0.8em]
{\large 题目、分步解析与数值核验}\\[0.5em]
{\small \textcolor{red!65!black}{审校稿：未经化学专家核准；禁止直接视为正式试卷或出版物}}
\end{center}
\vspace{1em}
\noindent 本册从仓库结构化 JSONL 自动生成。增补和修改应只编辑源题目数据，并重新执行统一编译。
\par 机械方程残差检查只能说明指定模型的内部数值一致，不保证模型现实可行性。第六辑另须完成跨辑模型去重审查。
\tableofcontents
\clearpage
\section*{第8辑 · 题目及解答}
\addcontentsline{toc}{section}{第8辑}
\subsection*{无机化学}
\addcontentsline{toc}{subsection}{第8辑 / 无机化学}
\Needspace{7\baselineskip}
\subsubsection*{UC-ADV8-I61 · 酸碱条件下配体与金属共同耗竭的非线性物料反演}
\noindent\textbf{题目}\par
M2+分析总浓度0.01200M，单酸配体HL总量0.01800M、Ka=10\textasciicircum{}-7，强缓冲pH8.000。仅L−与M2+形成ML+，形成常数K=2000L/mol，忽略其他反应。求未配位配体总量l、游离HL和L−、游离M与ML、结合率；给出完整双守恒与唯一根证明，不能将配体总量直接代入质量作用律。\par
\medskip\noindent\textbf{分步解答}
\begin{enumerate}
\item 令未配位配体池l=[HL]+[L−]，分布α=[L−]/l=Ka/(Ka+h)=0.909090909。
\item 金属守恒M=CM/(1+Kαl)，ML=CM−M。
\item 配体守恒Ltot=l+CMKαl/(1+Kαl)，左侧严格递增而唯一。
\item 数值l=0.00688740616M、M=0.000887406155M、ML=0.0111125938M。
\item HL=0.000626127832M、L−=0.00626127832M，结合率=0.926049487。
\item 核查上述每个物种的物料与电荷守恒，检查根的非负性、唯一性或不可能条件。
\item 独立重代平衡与守恒的数值残差：-1.927e-16、0.000e+0、0.000e+0。计算核对不替代真实材料与专家审校。
\end{enumerate}
\noindent\textbf{答案摘要}\par
答案1 = 0.00688740616；答案2 = 0.000626127832；答案3 = 0.00626127832；答案4 = 0.000887406155；答案5 = 0.0111125938\par
\noindent\textbf{模型前提：}
自由HL/L−酸碱平衡快、pH外部固定，理想单络合。\par
\noindent\textcolor{red!65!black}{[跨辑去重未完成 · 草稿]} \par
\medskip\hrule\medskip
\Needspace{7\baselineskip}
\subsubsection*{UC-ADV8-I62 · 四价施主掺杂与氧空位补偿符号的不可行性证明}
\noindent\textbf{题目}\par
假想A\_(1−x−y)D\_xE\_yO\_(1.5−δ)，A3+、D4+、E4+、O2−，x=0.1400、y=0.0600。限制氧空位δ≥0且不允许电子、阳离子空位、氧间隙。判断是否存在电中性解，并求形式氧缺陷量；若允许电子补偿但氧固定1.500，计算每化学式电子数、0.2500mol样品的电子mol及相应电量。不能把负氧空位解释成真实缺陷。\par
\medskip\noindent\textbf{分步解答}
\begin{enumerate}
\item 施主增加净正电量x+y=0.2每化学式。
\item 移除氧O2−再增加正电2δ，电中性要求x+y+2δ=0。
\item 形式解δ=−(x+y)/2=-0.1<0，违反氧空位δ≥0，故限制体系不存在物理解。
\item 若保持O1.5并允许电子，每化学式须引入x+y=0.2电子。
\item 0.2500mol样品电子量=0.05mol，对应电量=4824.26661C。
\item 核查上述每个物种的物料与电荷守恒，检查根的非负性、唯一性或不可能条件。
\item 独立重代平衡与守恒的数值残差：0.000e+0、0.000e+0、0.000e+0。计算核对不替代真实材料与专家审校。
\end{enumerate}
\noindent\textbf{答案摘要}\par
答案1 = -0.1；答案2 = 0.2；答案3 = 0.05；答案4 = 4824.26661\par
\noindent\textbf{模型前提：}
形式负δ仅用于不可行反证，不视为实在晶体。\par
\noindent\textcolor{red!65!black}{[跨辑去重未完成 · 草稿]} \par
\medskip\hrule\medskip
\Needspace{7\baselineskip}
\subsubsection*{UC-ADV8-I63 · 氧化态选择性质子化引起的Nernst条件电势偏移}
\noindent\textbf{题目}\par
一电子Ox+e−⇌Red，E°=0.5200V、T=298.15K，Ox总量0.002000M、Red总量0.001000M。仅Ox可质子化成HOx，常数KH=[HOx]/([Ox][H+])=10\textasciicircum{}4 L/mol，pH5.000由缓冲固定。求各氧化态物种、正确Nernst电势及误用总Ox的电势偏差，并检查电位对pH的方向。\par
\medskip\noindent\textbf{分步解答}
\begin{enumerate}
\item 游离氧化态分数f=1/(1+KHh)=0.909090909。
\item Ox游离=0.00181818182M，HOx=0.000181818182M，Red游离=Red总量。
\item 按游离物种应用E=E0+(RT/F)ln([Ox]/[Red])。
\item 正确E=0.535359974V，误算Etotal=0.537808739V。
\item 高质子化降低游离氧化态，误用总量会高估电位0.00244876434V。
\item 核查上述每个物种的物料与电荷守恒，检查根的非负性、唯一性或不可能条件。
\item 独立重代平衡与守恒的数值残差：0.000e+0、2.220e-16、3.990e-17。计算核对不替代真实材料与专家审校。
\end{enumerate}
\noindent\textbf{答案摘要}\par
答案1 = 0.00181818182；答案2 = 0.000181818182；答案3 = 0.535359974；答案4 = 0.00244876434\par
\noindent\textbf{模型前提：}
HOx仅假想氧化态质子物种，Red不质子化，忽略活度效应。\par
\noindent\textcolor{red!65!black}{[跨辑去重未完成 · 草稿]} \par
\medskip\hrule\medskip
\Needspace{7\baselineskip}
\subsubsection*{UC-ADV8-I64 · 两级络合的配分多项式与每金属配体平均占有数}
\noindent\textbf{题目}\par
假想金属M仅形成M、ML、ML2，K1=300L/mol、K2=50L/mol。M总浓度0.008000M，配体L总量0.01200M，忽略单配体其它酸碱形态。计算游离配体、三种金属形态以及平均每金属中心结合配体数；证明有限配体方程唯一，并指出不能仅由K2<K1判断微观反协同性。\par
\medskip\noindent\textbf{分步解答}
\begin{enumerate}
\item 结合配分多项式Z=1+K1l+K1K2l²。
\item 自由M=Mtot/Z，ML=K1lM、ML2=K1K2l²M。
\item 配体守恒Ltot=l+ML+2ML2，左侧对l严格增。
\item 求l=0.00544385406M，M=0.0025993526、ML=0.00424514886、ML2=0.00115549854M。
\item 平均配位数ν=(ML+2ML2)/Mtot=0.819518243，逐级常数含统计因素。
\item 核查上述每个物种的物料与电荷守恒，检查根的非负性、唯一性或不可能条件。
\item 独立重代平衡与守恒的数值残差：0.000e+0、0.000e+0、0.000e+0。计算核对不替代真实材料与专家审校。
\end{enumerate}
\noindent\textbf{答案摘要}\par
答案1 = 0.00544385406；答案2 = 0.0025993526；答案3 = 0.00424514886；答案4 = 0.00115549854；答案5 = 0.819518243\par
\noindent\textbf{模型前提：}
条件形成常数已吸收构型统计简并，M仅三种水相配位态。\par
\noindent\textcolor{red!65!black}{[跨辑去重未完成 · 草稿]} \par
\medskip\hrule\medskip
\Needspace{7\baselineskip}
\subsubsection*{UC-ADV8-I65 · 配位场两构型简并配分的谱学响应与热容耦合}
\noindent\textbf{题目}\par
某假想无机配位中心有低态L（能量0、简并度1、磁响应0.400）与高态H（能量高7.800kJ/mol、简并度5、磁响应1.400），310K两态独立达热平衡。求高态比例、热平均磁响应、相对能量和由两态重分布带来的附加热容，并解释为什么响应可线性加权而热容需要对概率求温度导数。\par
\medskip\noindent\textbf{分步解答}
\begin{enumerate}
\item 高态统计权重wH=5exp(−ΔE/RT)，配分Z=1+wH。
\item 高态概率pH=wH/Z=0.195168487，低态比例1−pH。
\item 平均磁响应s=1.4pH+0.4(1−pH)=0.595168487。
\item 平均相对能量U=ΔE pH=1522.3142J/mol。
\item 附加热容C=R(ΔE/RT)²pH(1−pH)=11.9604164J/(mol K)。
\item 核查上述每个物种的物料与电荷守恒，检查根的非负性、唯一性或不可能条件。
\item 独立重代平衡与守恒的数值残差：0.000e+0、0.000e+0、0.000e+0。计算核对不替代真实材料与专家审校。
\end{enumerate}
\noindent\textbf{答案摘要}\par
答案1 = 0.195168487；答案2 = 0.595168487；答案3 = 1522.3142；答案4 = 11.9604164\par
\noindent\textbf{模型前提：}
忽略其它磁耦合及振动热容，高低态响应不随温度变。\par
\noindent\textcolor{red!65!black}{[跨辑去重未完成 · 草稿]} \par
\medskip\hrule\medskip
\end{document}
